% ---------------------------------------------------------------------------
% Discussion section. Use with:  \input{discussion}
%
% All citation keys now match references.bib. The heatmap is referenced as
% Figure~0. Group-level values were read from
% Figures 1, 2 and 5 and are approximate: replace with exact values.
% ---------------------------------------------------------------------------

\section{Discussion}

\subsection{Principal findings}

In 59,570 patients with non-mild ischaemic stroke who arrived within six hours of onset at MT-performing hospitals in England and Wales, every reperfusion strategy we examined was associated with a higher chance of leaving hospital with little or no disability (mRS 0--2), and the size of that benefit depended strongly on when each treatment was given. Early combined treatment, with both IVT and MT delivered within 270 minutes of onset, produced the largest benefit by all three methods: about 46\% of treated patients against 24\% of their matched untreated controls were discharged with mRS 0--2, the S-learner estimated a log odds ratio of about 1.5 (odds ratio about 4.4), and the causal forest estimated an absolute gain of about 18 percentage points. Early MT alone (log odds ratio about 1.3; absolute gain about 14 percentage points) and early IVT alone (log odds ratio about 0.8; absolute gain about 10 percentage points) each gave substantial benefit. Benefit was consistently smaller when the same treatments were given later: late IVT alone and late MT alone gave roughly half to two thirds of the benefit, on the log odds scale, of their early counterparts, and patients who received early IVT but late MT gained little more than patients who received early IVT alone.

When onset-to-treatment times were modelled as continuous variables, the benefit of each treatment given alone declined approximately linearly. For IVT the estimated benefit fell from a log odds ratio of 1.19 extrapolated to the moment of onset to zero at 405 minutes (6.75 hours). For MT it fell from 1.54 at onset to zero at 596 minutes (9.9 hours). These correspond to a loss of about 0.18 and 0.16 log odds units per hour of delay respectively, so that each hour of delay multiplies the odds ratio for a good outcome by about 0.84 for IVT and 0.86 for MT. The joint analysis of both treatment times (Figure~\ref{fig:main_s_learner_matrix}) showed that these two clocks run together: for any given time to MT, benefit was greater the earlier IVT had been given, and for any given time to IVT, benefit was greater the earlier MT followed.

\subsection{What this study adds: time to IVT and time to MT considered together}

To our knowledge, this is the first study to describe how the benefit of reperfusion treatment varies across the joint distribution of onset-to-IVT and onset-to-MT times in a national population not selected by advanced imaging. The existing evidence treats the two clocks separately. The alteplase trials relate benefit to time to IVT in populations for whom MT was not available \cite{emberson_effect_2014}. The HERMES analyses relate benefit to time to arterial puncture or reperfusion in patients most of whom had already received alteplase, without modelling when the alteplase was given \cite{goyal2016,saver2016}. The trials of MT with or without prior IVT randomised IVT at a single moment in patients who were already at an MT centre, so that the interval between IVT and MT was short and varied little \cite{kaesmacher2024,majoie2023}. None of these designs can show what happens when the two times vary independently. In routine practice they do: many patients are treated with IVT at one hospital and undergo MT at another, one or several hours later.

Figure~\ref{fig:main_s_learner_matrix} describes this joint dependence, and three patterns emerge. First, for a fixed time to MT, benefit rises the earlier IVT was given. For MT at 240--300 minutes, the estimated log odds ratio was 0.84 without IVT, 0.81--0.82 when IVT was given at 180--300 minutes, 0.99 when it was given at 120--180 minutes, and 1.29 when it was given at 60--120 minutes. IVT therefore adds substantially to MT when given within about two hours of onset (an additional 0.3--0.6 log odds units across MT times of 3--6 hours, equivalent to odds ratios of about 1.4--1.8), adds less at 2--3 hours, and adds little when first given after three hours and followed promptly by MT.

Second, for a fixed time to IVT, benefit falls the later MT follows, and converges on the benefit of IVT alone. For IVT given at 60--120 minutes, the log odds ratio was 1.70 with MT at 120--180 minutes, 1.37 at 180--240 minutes, 1.29 at 240--300 minutes and 1.11 at 300--360 minutes, against 0.98 for IVT alone in the same time band. The same pattern was seen for IVT at 120--180 minutes (1.78 falling to 0.70 by 360--420 minutes, against 0.73 for IVT alone). In this population the additional value of MT over early IVT was therefore largely lost by six to seven hours after onset. The categorical analysis says the same thing more simply: patients given early IVT and late MT had an estimated benefit (log odds ratio about 0.83) barely above that of early IVT alone (about 0.77).

Third, when IVT was given later, it appeared to add more when MT was delayed further still. IVT at 180--240 minutes added nothing to MT performed in the same or the following hour (1.03 against 1.03; 0.82 against 0.84), but was associated with an additional 0.2--0.3 log odds units when MT did not take place until 300--480 minutes. One interpretation is that IVT matters most when the wait for MT is long, because the drug has more time to act; recanalisation before MT occurs in around one fifth of IVT-treated patients referred for MT and is more frequent in those transferred between hospitals than in those treated on one site \cite{seners2016,seners2018}. Cell sizes here are small, and we regard this observation as hypothesis-generating.

The benefits of the two treatments were less than additive. The log odds ratio for early combined treatment (about 1.5) was well below the sum of those for early IVT alone and early MT alone (about 0.8 and 1.3). This is expected of two treatments that act through the same mechanism, reperfusion of the same tissue, and it means that neither the IVT trials nor the MT trials can be used on their own to predict the value of a combined strategy at a given pair of treatment times.

\subsection{Comparison with randomised trials of IVT}

Our estimates for IVT alone agree closely with the randomised evidence in absolute terms and in their dependence on time. In the individual patient data meta-analysis of Emberson et al., alteplase increased the proportion of patients with an excellent outcome by about 10 percentage points when given within 3 hours and about 5 percentage points at 3--4.5 hours \cite{emberson_effect_2014}. The causal forest estimated an absolute gain in discharge mRS 0--2 of about 10 percentage points for IVT within 4.5 hours and about 9 percentage points after 4.5 hours, and the matched analysis gave about 11 and 7 percentage points. Emberson et al.\ estimated that the benefit of alteplase reached zero at 6.3 hours after onset (95\% CI 5.0--13.8) \cite{emberson_effect_2014}; our linear model placed this point at 6.75 hours. The agreement is notable given that the two estimates come from different designs, eras and outcome definitions, although ours is extrapolated beyond the six hours within which IVT was given in our data.

Our relative effects were larger than those of the trials. The S-learner odds ratio for early IVT alone was about 2.2, against 1.75 within 3 hours and 1.26 at 3--4.5 hours in the meta-analysis \cite{emberson_effect_2014}. Several differences in design may contribute. We restricted the cohort to patients with NIHSS above 5 or a recorded large vessel occlusion, we measured mRS 0--2 at discharge and not mRS 0--1 at 3--6 months, and door-to-needle times and stroke unit care have improved since the trials. Residual confounding is also likely to act in this direction: patients who arrive early with a disabling stroke and are not treated with IVT usually have a reason, such as anticoagulation, recent surgery, frailty or a clinical judgement of futility, and not all of these reasons are captured in our covariates.

The rate at which benefit was lost with delay is consistent with earlier registry work. In 58,353 patients treated with alteplase in the Get With The Guidelines-Stroke registry, each 15 minutes of earlier treatment was associated with an odds ratio of 1.04 for independent ambulation at discharge \cite{saver2013}, equivalent to about 0.16 log odds units per hour. Our estimate for the decline in treatment effect on discharge mRS 0--2 was about 0.18 log odds units per hour. The registry analysis could compare only treated patients with one another; by emulating an untreated comparison group we were able to estimate the treatment effect itself at each time, and so to locate the time at which benefit is exhausted. Time dependence of alteplase has also been examined specifically in large vessel occlusion: in the control arms of the HERMES trials, each hour of delay was associated with an odds ratio of 0.82 (95\% CI 0.66--1.03) for functional independence \cite{goyal2019}.

We did not examine intracranial haemorrhage. The trials show that alteplase increases the risk of fatal intracranial haemorrhage by about 2 percentage points irrespective of treatment delay \cite{emberson_effect_2014,whiteley2016}, which is why the net benefit of late IVT is small. Our outcome, disability at discharge, includes the consequences of haemorrhage that occurred during the admission, so the estimates presented are net of early harm, but we cannot separate the two.

\subsection{Comparison with randomised trials of MT}

The time dependence we estimated for MT is very close to that reported from the trials. In the HERMES time-to-treatment analysis, each hour of delay to reperfusion was associated with a common odds ratio of 0.84 for less disability and an odds ratio of 0.81 for functional independence \cite{saver2016}. In MR CLEAN, the absolute treatment effect fell by about 6 percentage points per hour of delay to reperfusion \cite{fransen_time_2016}. Our estimate for MT alone, a loss of about 0.16 log odds units per hour, corresponds to an odds ratio of 0.86 per hour. In HERMES the benefit of MT was no longer statistically significant when arterial puncture was later than 7 hours 18 minutes after onset, although the point estimate still favoured treatment at 8 hours (common odds ratio 1.57) \cite{saver2016}. Our point estimate of the time at which benefit reached zero, 9.9 hours, is compatible with this, and in Figure~\ref{fig:main_s_learner_matrix} the estimated benefit of MT alone was close to zero in each time band beyond 9 hours.

The size of the benefit requires more careful comparison, because the comparator differs. The HERMES trials compared MT plus standard care with standard care alone in patients with a proven anterior circulation large vessel occlusion, most of whom received alteplase in both arms; functional independence at 90 days was achieved by 46.0\% and 26.5\% \cite{goyal2016}. Our treated groups are compared with patients who received no reperfusion treatment. On that basis, early MT alone was associated with an absolute gain of about 14--17 percentage points and early combined treatment with a gain of about 18--22 percentage points, and the proportions with a good outcome in the early combined group and its controls (about 46\% and 24\%) are similar to those in the two arms of HERMES. The trial contrast, MT added to IVT, is best approximated in our data by the difference between combined treatment and IVT alone at the same time to IVT. For IVT given at 60--180 minutes, MT within 180 minutes was associated with an additional 0.7--1.0 log odds units (odds ratios of about 2.0--2.9), which brackets the HERMES common odds ratio of 2.49 \cite{goyal2016} and the 2.79 estimated at 3 hours \cite{saver2016}. At later MT times the increment in our data fell faster than in the trials: to about 0.3--0.4 log odds units at 3--5 hours and about 0.1 at 5--6 hours, against a common odds ratio of 1.98 at 6 hours in HERMES \cite{saver2016}.

There are several reasons to expect a smaller increment at later times in routine practice than in trials. Our IVT-only group includes patients without a large vessel occlusion and patients whose occlusion recanalised with IVT, so its outcomes are better than those of a trial control arm in which every patient had a persisting occlusion. Conversely, patients who proceed to MT after IVT are by definition those in whom IVT has not yet succeeded. Several of the HERMES trials selected patients using collateral or perfusion imaging, which enriches for slowly progressing infarcts in whom late reperfusion still helps \cite{goyal2016,saver2016}. Finally, disability at discharge may understate the benefit of MT, since patients with severe strokes continue to recover over the following months.

The patients who received MT without IVT were a distinct group: 39--53\% were taking an anticoagulant for atrial fibrillation, against 3--4\% of those who received both treatments. Their results therefore describe MT in patients who were ineligible for IVT, and should not be read as a comparison of direct MT with bridging in patients eligible for both. The finding that MT was effective in this group is consistent with HERMES, in which no heterogeneity of treatment effect was found across prespecified subgroups \cite{goyal2016}, and with MR CLEAN, in which the effect of treatment did not differ between patients who had and had not received alteplase \cite{mulder2016}.

\subsection{Comparison with trials of MT with and without prior IVT}

The randomised evidence closest to our joint analysis comes from the six trials that compared MT alone with IVT followed by MT in patients presenting directly to MT centres: DIRECT-MT, DEVT, SKIP, MR CLEAN-NO IV, SWIFT DIRECT and DIRECT-SAFE \cite{yang2020,zi2021,suzuki2021,lecouffe2021,fischer2022,mitchell2022}. Individually these trials were inconsistent. Pooled by the IRIS collaboration (2,313 patients), they did not establish non-inferiority of MT alone: the adjusted common odds ratio for MT alone against combined treatment was 0.89 (95\% CI 0.76--1.04) \cite{majoie2023}. The average effect of adding IVT immediately before MT is therefore small.

A subsequent IRIS analysis showed that this average conceals a strong dependence on time \cite{kaesmacher2024}. The benefit of IVT plus MT over MT alone fell with time from onset to expected IVT (ratio of adjusted common odds ratios 0.84 per hour): the adjusted common odds ratio was 1.49 at 1 hour, 1.25 at 2 hours and 1.04 at 3 hours, with absolute differences in functional independence of 9, 5 and 1 percentage points. The benefit was no longer statistically significant after 2 hours 20 minutes, and the point estimate crossed the null at 3 hours 14 minutes \cite{kaesmacher2024}.

Our observational estimates reproduce this pattern closely. In the IRIS trials MT followed IVT quickly (median onset to arterial puncture was 180 minutes, about half an hour after the median expected time of IVT) \cite{kaesmacher2024}, so the relevant cells of Figure~\ref{fig:main_s_learner_matrix} are those in which MT took place in the same hour as IVT or the next. In those cells, the estimated benefit with IVT at 60--180 minutes exceeded that of MT alone at the same time by about 0.1--0.3 log odds units (odds ratios of about 1.1--1.35), and with IVT after 180 minutes the difference was nil (1.03 against 1.03 at 180--240 minutes; 0.81 against 0.84 at 240--300 minutes). That two entirely different designs, one randomised and one emulated from registry data, locate the loss of added benefit at about three hours strengthens both findings.

Our data extend the trial evidence in two directions that the trials could not address. First, the IRIS authors caution that their results apply only to patients who present directly to an MT centre \cite{kaesmacher2024}. In the UK, IVT at a local unit followed by secondary transfer is the predominant model of MT provision \cite{allen2022}, and in our data MT frequently followed IVT by two hours or more. Figure~\ref{fig:main_s_learner_matrix} covers these patients, and suggests that for them early IVT is worth considerably more than the trials of direct presenters imply (an additional 0.3--0.6 log odds units when IVT is given at 60--120 minutes and MT follows at 3--6 hours). Second, the trials varied only whether IVT was given, at one point in time; we describe how benefit changes as each treatment time moves while the other is held fixed.

More recent trials are consistent with our findings. In BRIDGE-TNK, in patients presenting within 4.5 hours, tenecteplase before MT increased functional independence at 90 days from 44.1\% to 52.9\% \cite{qiu2025}. In TNK-PLUS, in patients presenting 4.5--24 hours after onset, tenecteplase before MT made no difference (44.2\% against 43.2\%) \cite{tnkplus2026}. A secondary analysis of SWIFT DIRECT alone did not find clear evidence that the effect of bridging depended on onset-to-needle time, although it was not powered to do so \cite{meinel2022}.

\subsection{Comparison with registry and other observational evidence}

Large registries have confirmed, among treated patients, that the time dependence seen in trials holds in routine care. In the Get With The Guidelines-Stroke registry (6,756 patients, onset to puncture within 8 hours), earlier MT was associated with better outcomes at discharge, and the relationship was non-linear, with a steeper gradient between 30 and 270 minutes than between 271 and 480 minutes \cite{jahan2019}. Our continuous estimates for MT alone show the same shape, falling steeply over the first four hours and more slowly thereafter, and it is of interest that the inflection identified in that registry coincides with the 270-minute threshold used to define early and late treatment here. In the MR CLEAN Registry (1,488 patients), each hour of delay to the start of MT was associated with an adjusted common odds ratio of 0.83 for functional outcome \cite{mulder2018}, almost identical to our 0.86.

Observational comparisons of MT with and without prior IVT have generally favoured combined treatment. In the MR CLEAN Registry, IVT before MT was associated with better functional outcome (adjusted common odds ratio 1.47) \cite{chalos2019}, and in a recent analysis of 7,118 patients from 46 centres it was associated with functional independence whether MT took place within or after 4.5 hours (adjusted odds ratios 1.56 and 1.27), with no significant interaction with the MT time window \cite{chacon2026}. These analyses classify patients by whether IVT was given and by when MT took place, but not by when IVT was given. Our results suggest that this third quantity is the important one: the time of MT alone does not determine whether IVT adds benefit, whereas the time of IVT, and the interval that follows it, do. Among patients who received both treatments in a Berlin registry, a shorter interval between IVT and MT was associated with better outcome at three months \cite{wagner2022}, which corresponds to our finding that, for a fixed time to IVT, benefit falls as MT is delayed.

These registry studies share a limitation that our design addresses in part. They compare treated patients with other treated patients, so they can estimate how outcome varies with time or with co-treatment, but not the effect of treatment against no treatment at each time. By constructing matched or model-based untreated comparators from the same national registry, we estimated treatment effects directly, at the cost of stronger assumptions about confounding, discussed below.

\subsection{Late treatment and the late-window trials}

Our finding that the benefit of MT declines to nothing by about ten hours appears to conflict with the trials of late MT, but the populations are different. DAWN and DEFUSE 3 enrolled patients 6--24 hours after they were last known to be well who had a small infarct core relative to their clinical deficit or perfusion lesion \cite{nogueira2018,albers2018}. In the AURORA pooled analysis of patients randomised more than six hours after onset, functional independence was achieved by 45.9\% with MT and 19.3\% without (adjusted common odds ratio 2.54), an effect as large as that seen in the early window \cite{jovin2022}. MR CLEAN-LATE found benefit at 6--24 hours in patients selected only by the presence of collateral flow on CT angiography \cite{olthuis2023}, and observational data suggest that patients selected in the late window by non-contrast CT have outcomes similar to those selected by perfusion imaging \cite{nguyen2022}.

Patients in those trials presented late and were chosen because imaging showed that salvageable tissue remained; they are, by selection, patients whose infarcts progress slowly. The patients we studied all arrived within six hours of a known time of onset, and those treated late were, we presume, treated late mainly because of delay, not because imaging had identified them as late-window candidates. Our estimates therefore describe what delay costs the general population of early presenters, in whom fast and slow progressors are mixed. They do not describe the effect of MT in imaging-selected late presenters, and should not be used to argue against late-window treatment. The two bodies of evidence are complementary: the late-window trials show who can still benefit when presentation is late, and our data show how much is lost, on average, when treatment of an early presenter is delayed.

Late MT was nonetheless associated with a clear benefit in our data. MT alone after 270 minutes was associated with an absolute gain in discharge mRS 0--2 of about 7--8 percentage points, and all three methods found this to be significantly greater than zero. The number of patients treated beyond nine hours was small, and estimates in that range are imprecise. The group that received both treatments late contained only 141 patients, and its wide confidence intervals do not allow any conclusion; the randomised evidence for that combination is the TNK-PLUS trial, which found no benefit of IVT before MT at 4.5--24 hours \cite{tnkplus2026}.

\subsection{Agreement between methods, and who benefits}

The three causal methods gave substantially the same results. Matching, the S-learner and the causal forest make different modelling assumptions: matching compares like with like on a small set of prognostic features without a model of outcome, the S-learner relies on a flexible outcome model, and the causal forest combines an outcome model with a model of treatment assignment and is consistent if either is correctly specified. All three ranked the strategies in the same order, and where they estimate the same quantity the values were close: the absolute gain in discharge mRS 0--2 from matching and from the causal forest was about 11 and 10 percentage points for early IVT alone, 17 and 14 for early MT alone, and 22 and 18 for early combined treatment. Agreement between methods does not remove bias from confounders that none of them observe, but it makes it unlikely that the findings are an artefact of one model.

The agreement with randomised trials, where trials exist, serves as a further check. The emulation recovered the absolute benefit of IVT and the time at which it is exhausted \cite{emberson_effect_2014}, the hourly decay of MT benefit \cite{fransen_time_2016,saver2016}, and the time dependence of the added value of IVT before MT \cite{kaesmacher2024}. Benchmarking an emulated target trial \cite{hernan2016} against known trial results is a recognised means of assessing whether observational data can support causal conclusions, and it gives some confidence in the estimates for which no trial exists, namely those for specific combinations of the two treatment times.

Our estimates are of the effect of treatment in those who received it, and the treated groups differed from one another. Patients who received MT had more severe strokes and less prior disability than those who received IVT alone, and those who received MT without IVT were far more likely to have atrial fibrillation and to be anticoagulated. Differences in estimated benefit between groups therefore reflect differences in the patients as well as in the treatments and their timing. The secondary model of estimated benefit indicated which characteristics matter most. Benefit increased with stroke severity and decreased with pre-stroke disability across mRS 0--4, as would be expected for an outcome defined as mRS 0--2, and the admitting stroke team was the single most influential feature. Variation between teams of this size, after accounting for patient characteristics and treatment times, suggests that how treatment is delivered, including patient selection, procedural success and post-procedural care, influences how much patients gain. It may also reflect differences between teams in case mix that we did not measure and in how discharge disability is assessed, and it deserves separate study.

\subsection{Strengths and limitations}

The study has several strengths. It uses a national registry covering stroke admissions across England and Wales over six years, giving a cohort around nine times the size of the pooled IVT trials and far larger than the pooled MT trials. Patients were not selected by advanced imaging and were not subject to trial exclusion criteria for age or prior disability, so the results describe the patients who are actually treated. Three causal methods with different assumptions were applied within an explicit target trial framework, treatment times were analysed both in categories and continuously, and all code is openly available.

The principal limitation is that treatment was not randomised. We adjusted for the main determinants of outcome and of treatment that the registry records, but several important ones were not available. We had no imaging data: occlusion site, early ischaemic change and collateral status all influence both the decision to treat and the outcome. A large vessel occlusion is confirmed in every patient who undergoes MT but not in their untreated comparators, who were matched on stroke severity; if comparators include patients without an occlusion, who have a better prognosis, the benefit of MT will be underestimated. In the opposite direction, patients who arrive early with a severe stroke and receive no reperfusion treatment have often been judged unsuitable for reasons such as frailty or extensive established infarction, which would lead to overestimation. We cannot say which bias dominates.

The treatment groups are defined by what happened after arrival, which introduces selection that a trial would avoid. Patients whose occlusion recanalised after IVT did not proceed to MT and are counted in the IVT-only group; the combined-treatment groups therefore contain those in whom IVT had not yet succeeded. This favours IVT alone and disadvantages combined treatment, particularly late MT after early IVT. Patients must also survive, and remain candidates, until the time of a late treatment, whereas untreated comparators are followed from arrival.

Outcome was disability at discharge, assessed by the treating team without blinding, and not at 90 days as in the trials. Discharge mRS predicts 90-day mRS, but imperfectly, and scores continue to change between 30 and 90 days after discharge \cite{elhabr2021}. If patients with severe strokes treated by MT recover more after discharge than untreated patients do, our estimates of its benefit will be conservative. Length of stay differed between treated and untreated groups, so disability was not measured at the same time after stroke in each.

Several further points should be borne in mind. Time of onset was known for all patients included, so the findings do not extend to wake-up stroke or stroke of unknown onset. The threshold of 270 minutes for early and late treatment reflects the licensed window for alteplase and is otherwise arbitrary, although the continuous analyses do not depend on it. Some combinations of treatment times were rare, and the times at which benefit reached zero are extrapolations of a linear fit. The confidence intervals for the S-learner describe the precision of the mean predicted effect and do not include uncertainty in the model itself, so they are narrower than the true uncertainty. We did not distinguish alteplase from tenecteplase, nor examine symptomatic intracranial haemorrhage or mortality separately. Finally, the cohort was limited to hospitals that performed at least ten MT procedures and to patients with NIHSS above 5 or a recorded large vessel occlusion, and the results may not apply to milder strokes or to other health systems.

\subsection{Implications for practice, services and research}

For clinicians, the results support giving IVT as early as possible to eligible patients who are also candidates for MT, and particularly to those who face a wait or a transfer before MT can begin. They give no support to withholding IVT in order to proceed to MT when IVT can be given early. They also suggest, in agreement with the randomised evidence \cite{kaesmacher2024,tnkplus2026}, that IVT first given more than about three hours after onset adds little when MT follows promptly. That observation concerns average benefit in an observational analysis, and is not a basis for changing individual treatment decisions.

For those who plan services, the central finding is that most of the additional value of MT over early IVT was lost when MT was delayed beyond about six hours. The national MT rate in England was 3.9\% in the 2024 national audit report, and median door-in-door-out times at referring hospitals of two to three hours have been reported \cite{nagaratnam2025}. Our results indicate that delay of this order is not only a barrier to access but removes much of the benefit for those who are eventually treated. Shortening the interval between IVT and MT, whether by faster referral and transfer or by bringing MT closer to patients, is likely to yield substantial gains.

The joint estimates also bear on the choice between conveying patients directly to an MT centre and treating first at the nearest unit. Bypassing the local unit brings MT forward and pushes IVT back. In the RACECAT trial, direct transport to an MT centre in non-urban Catalonia reduced the proportion of patients treated with IVT, increased the proportion receiving MT, and made no difference to disability at 90 days \cite{perezdelaossa2022}. Figure~\ref{fig:main_s_learner_matrix} shows why such a result is plausible: the gain from earlier MT and the loss from later IVT are of broadly similar size over the relevant range, so the net effect depends on the local travel times. A grid of estimated benefit by both treatment times provides the quantity that models of service configuration require but have so far had to assemble from separate IVT and MT trials \cite{allen2022}.

Further work should test these findings where our data are weakest. Linkage to imaging would allow comparators to be restricted to patients with a confirmed large vessel occlusion. Outcomes at 90 days or later, and safety outcomes, should be examined. The apparent dependence of the value of IVT on the interval before MT needs confirmation in a dataset that records transfers and recanalisation before MT. As tenecteplase replaces alteplase, the relationship should be re-estimated, since the two drugs may differ in how quickly they act. The large variation in benefit between stroke teams also needs explanation.

\subsection{Conclusion}

In patients with non-mild ischaemic stroke who reach hospital within six hours, IVT and MT are each effective in routine practice in England and Wales, with benefits and time dependence that closely match those of the randomised trials. Considering the two treatment times together shows what the trials, which studied each in isolation, could not: early IVT adds substantially to MT, especially when MT is delayed; late IVT adds little when MT follows promptly; and MT adds little to early IVT once it is delayed beyond about six hours. The greatest benefit comes from giving both treatments early, and the organisation of stroke services should aim to minimise both times, not either one alone.